\chapter*{TÓM TẮT}
\addcontentsline{toc}{chapter}{TÓM TẮT}

Người dân thường hỏi về luật giao thông bằng lời nói thường ngày, có khi gõ không dấu hoặc viết tắt, trong khi câu trả lời đúng nằm rải rác trong nhiều văn bản quy phạm pháp luật có hàng nghìn Điều, Khoản, Điểm và nhiều lần sửa đổi với ngày hiệu lực khác nhau. Một hệ thống RAG thông thường cắt văn bản thành đoạn cố định và tìm theo vector thường làm mất cấu trúc pháp lý, tìm lệch chủ đề với câu khẩu ngữ, và không biết điều khoản nào đã bị sửa đổi.

Đồ án xây dựng một hệ thống truy hồi thông tin và hỏi đáp cho lĩnh vực luật giao thông đường bộ Việt Nam trên 14 văn bản còn hiệu lực. Luồng nạp chuyển văn bản nguồn thành cây Điều--Khoản--Điểm dựa trên quy ước soạn thảo, giải các liên kết trích dẫn trên cây phân tích và gắn ghi chú sửa đổi đọc từ chính các câu lệnh sửa đổi. Tầng tìm kiếm kết hợp vector Qwen3-Embedding với BM25 trên âm tiết không dấu bằng Reciprocal Rank Fusion, cộng thêm tìm theo cụm từ nguyên văn và lọc theo ngày vi phạm. Một agent dựa trên mô hình ngôn ngữ lớn gọi các công cụ tra cứu qua giao thức MCP để trả lời nhiều bước, có kiểm tra số liệu bám nguồn.

Hệ thống được đánh giá hai tầng, tách tập \texttt{dev} và tập \texttt{holdout}. Trên bộ 715 câu, truy hồi kết hợp tìm thấy đáp án ở 93,4\% câu \texttt{dev} và 97,6\% câu \texttt{holdout}; riêng việc thay \texttt{ts\_rank} bằng BM25 nâng Hit@1 của nhánh từ khoá từ 13,7\% lên 57,8\%. Trên bộ 81 câu theo lời hỏi thật, agent trả lời đúng toàn phần 97,2\% câu trong phạm vi, không từ chối nhầm và không bịa mức tiền, đồng thời từ chối đúng 8/10 câu ngoài phạm vi.

\textbf{Từ khoá:} truy hồi thông tin, RAG, văn bản quy phạm pháp luật, luật giao thông đường bộ, truy hồi kết hợp, BM25, agent, MCP.

\section*{ABSTRACT}

\begin{sloppypar}
People ask about traffic law in everyday language, often without Vietnamese diacritics, while the answer is spread across statutes with thousands of articles, clauses and points, amended at different dates. Fixed-size chunking with vector search loses the legal structure, drifts off topic on colloquial questions and ignores amendments.

This project builds a retrieval and question-answering system over 14 Vietnamese road traffic statutes. Ingestion turns source documents into an article--clause--point tree based on legislative drafting conventions, resolves cross-references on the parse tree and attaches amendment notes read from the amending instructions themselves. Search fuses Qwen3-Embedding vectors with BM25 over unaccented syllables using Reciprocal Rank Fusion, adds verbatim phrase matching and filters by the date of the violation. An LLM agent calls the lookup tools through the Model Context Protocol and checks that every figure it states appears in its sources.

On 715 retrieval cases, hybrid search finds the expected provision for 93.4\% of development and 97.6\% of held-out cases; replacing \texttt{ts\_rank} with BM25 alone raises keyword Hit@1 from 13.7\% to 57.8\%. On 81 user-style questions, the agent answers 97.2\% of in-scope questions fully correctly, with no wrong refusals and no invented fines, and refuses 8 of 10 out-of-scope questions.

\textbf{Keywords:} information retrieval, retrieval-augmented generation, legal documents, road traffic law, hybrid search, BM25, LLM agent, MCP.
\end{sloppypar}
